% Generated by code/assemble_manuscript.py; edit source fragments.

\section{B. Counterexamples and necessary fibre defects}

\subsection{B.1. The nodal test family and every gluing scalar}

% B1 from part_B.tex

\begin{statement}{B1}{Proved}{A complete nodal-family computation}
Let $0<\varepsilon<4/27$, let $\Delta=\{|t|<\varepsilon\}$, and define
\[
 \mathcal C=\{y^2z=x^2(x+z)+tz^3\}\subset\PP^2\times\Delta,
 \qquad X=\mathcal C\times\PP^1.
\]
The projection $X\to\Delta$ is proper, its total space is smooth, and its central fibre is $S=C_0\times\PP^1$, where $C_0$ is an irreducible nodal cubic. For each $\lambda\in\C^*$ let $A_\lambda$ be the line bundle whose pullback to the normalization $\nu:\PP^1_u\to C_0$ is trivial and whose sections in this trivialization satisfy $h(1)=\lambda h(-1)$. Put $A=\operatorname{pr}_1^*A_\lambda$ and $d_\lambda=1$ for $\lambda=1$, and $d_\lambda=0$ otherwise. In the order $q=0,1,2$, the coherent cohomology dimensions are
\[
\begin{array}{c|c}
\text{sheaf on }S & (h^0,h^1,h^2)\\\hline
 N^*_{S/X}\otimes A &(d_\lambda,d_\lambda,0)\\
 K_S\otimes A &(1,0,0)\\
 T_S\otimes A &(1,0,0)\\
 \widetilde\Omega^1_S\otimes A &(0,1+d_\lambda,d_\lambda)\\
 \Omega^1_S\otimes A &(1,1+d_\lambda,d_\lambda)\\
 \Omega^1_X|_S\otimes A &(1,1+d_\lambda,d_\lambda).
\end{array}                                                    \tag{B.1}
\]
Here $T_S$ denotes the torsion of the cotangent sheaf, and $K_S$ its inverse image in the ambient cotangent sheaf. In particular the two proposed vanishings hold exactly for $\lambda\ne1$, and the ambient space of sections has dimension one for every $\lambda\in\C^*$. After shrinking $\Delta$, every $A_\lambda$ extends to a line bundle on $\mathcal C$ of the form $\OO_{\mathcal C}(P-Q)$ for holomorphic sections $P,Q$.
\end{statement}

\begin{proof}
In the affine chart $z=1$, the defining equation is $F=y^2-x^3-x^2-t=0$. Its partial derivative in $t$ is $-1$, so the total space is smooth there. Its only point with $z=0$ is $[0:1:0]$. In the chart $y=1$ the equation is $z-x^3-x^2z-tz^3=0$, whose $z$ derivative is one at this point. Fibre singularities in the affine chart satisfy $y=0$ and $x(3x+2)=0$; they occur at $t=0,-4/27$. Our disc contains only the first. The quadratic term $y^2-x^2$ at the origin has two distinct linear factors, and its smoothing parameter has nonzero differential. The holomorphic Morse lemma with parameter therefore supplies local coordinates with $t=vw$. Properness follows from the closed inclusion into $\PP^2\times\Delta$.

The normalization is
\[
 u\longmapsto(x,y)=(u^2-1,u(u^2-1)),
\]
completed by the point at infinity. The two points above the node are $1,-1$. The descent sequence, an exact sequence of coherent $\OO_{C_0}$-modules, is
\[
0\longrightarrow A_\lambda\longrightarrow\nu_*\OO_{\PP^1}
 \xrightarrow{\ h\mapsto h(1)-\lambda h(-1)\ }\C_{\rm node}
 \longrightarrow0.                                      \tag{B.2}
\]
Its cohomology is the two-term complex $\C\xrightarrow{1-\lambda}\C$. Thus $h^0(A_\lambda)=h^1(A_\lambda)=d_\lambda$.

Write $V=\Omega^1_{\mathcal C}|_{C_0}$. The normalization map induces a surjective $\OO_{C_0}$-linear morphism $\Omega^1_{C_0}\to\nu_*\Omega^1_{\PP^1}$; at a node $vw=0$, it sends $dv$ to $(dv,0)$ and $dw$ to $(0,dw)$. Its kernel is the length-one module generated by $v\,dw=-w\,dv$. Consequently
\[
 T_{C_0}=\C_{\rm node},\qquad
 \widetilde\Omega^1_{C_0}=\nu_*\OO_{\PP^1}(-2).
\]
The kernel of $V\to\nu_*\Omega^1_{\PP^1}$ is $\nu_*\OO_{\PP^1}$. Locally its generators are $w\,dv$ and $v\,dw$ on the two branches; the identification is the normal-crossing conormal identification of Part A with the globally fixed conormal $dt$. Tensoring by $A_\lambda$ preserves exactness, and finite pushforward and the projection formula identify the two resulting normalization terms with $\nu_*\OO$ and $\nu_*\OO(-2)$. Hence
\[
 (h^0,h^1)(C_0,V\otimes A_\lambda)=(1,1),\qquad
 (h^0,h^1)(C_0,\Omega^1_{C_0}\otimes A_\lambda)=(1,1).
\]
The product morphisms give the direct decompositions
\[
 \Omega^1_X|_S\otimes A
 =(V\otimes A_\lambda)\boxtimes\OO
 \ \oplus\ A_\lambda\boxtimes\OO(-2),
\]
and the analogous decomposition for $\Omega^1_S$. The torsion-free quotient is obtained by replacing the first factor by $\nu_*\OO(-2)$. The torsion is $\C_{\rm node}\boxtimes\OO$, its ambient inverse image is $\nu_*\OO\boxtimes\OO$, and $N^*_{S/X}=\OO_S$. The coherent Kunneth isomorphism and $H^0(\PP^1,\OO)=\C$, $H^1(\PP^1,\OO(-2))=\C$ now give every entry of (B.1). All other cohomology of these two line bundles on $\PP^1$ vanishes.

We give the extension explicitly. Fix $Q(0)$ at $u=0$, and take $P(0)$ at $u=a=(1-\lambda)/(1+\lambda)$; for $\lambda=-1$ take $a=\infty$. For $\lambda\ne0$ both points are in the smooth locus. With $b=0$, the rational function $r(u)=(u-a)/(u-b)$ has divisor $P(0)-Q(0)$, and a local section in its pullback trivialization descends precisely with scalar
\[
 \frac{r(1)}{r(-1)}
 =\frac{(1-a)(1+b)}{(1-b)(1+a)}=\lambda.
\]
For $a=\infty$ the same expression is its limit, namely $-1$. The projection $\mathcal C\to\Delta$ is a submersion at these points, so the holomorphic implicit-function theorem extends them to sections on a smaller disc. The point at infinity itself is a constant section. For $\lambda=1$ take $P=Q$ identically. Sections through the smooth fibre locus are effective Cartier divisors in the smooth surface $\mathcal C$; their difference therefore defines the claimed extension. Pullback along $X\to\mathcal C$ extends $A$. This proves the claim also for all roots of unity, including the identity.
\end{proof}

The calculation recovers the local example in the supplied correction note, Lemma 4.3 and Theorem 4.4, with the gluing convention fixed by (B.2). The restriction on roots of unity in Lead 2 is unnecessary: the only exceptional scalar for the conormal vanishing is $1$. The symbolic computation and the exact gluing matrices are recorded in \texttt{code/B\_examples.py} and its output.


\subsection{B.2. A compact counterexample to the ambient inference}

% B2 from part_B.tex

\begin{statement}{B2}{Proved}{A projective example with twelve length-one contributions}
Let $h:E\to\PP^1_t$ be the Weierstrass surface defined on the finite chart by
\[
 y^2z=x^3+(t^4+1)xz^2+t^6z^3,                          \tag{B.3}
\]
and at $s=1/t$ by $x_\infty=s^2x$, $y_\infty=s^3y$, $z_\infty=z$. Let $O=[0:1:0]$ and $P=(0,t^3,1)$ be its two sections and put
\[
 M_E=\OO_E(P-O),\quad Y=E\times\PP^1,\quad
 g=h\circ\operatorname{pr}_E,\quad
 L=K_Y\otimes\operatorname{pr}_E^*M_E^{-1},\quad
 \mathcal E=T_Y\otimes L.
\]
The surface $E$ is smooth and projective. Its singular fibres are twelve irreducible nodal cubics, at the distinct roots $t_i$ of
\[
 D(t)=31t^{12}+12t^8+12t^4+4.
\]
For every such fibre $S_i=g^{-1}(t_i)$, with $A=(K_Y\otimes L^{-1})|_{S_i}$, we have
\[
H^0(S_i,N^*_{S_i/Y}\otimes A)=0,
\qquad H^0(S_i,\widetilde\Omega^1_{S_i}\otimes A)=0,
\]
\[
 (h^0,h^1,h^2)(S_i,\mathcal E|_{S_i})=(0,1,1).
\]
As coherent sheaves on the base, with noncanonical identifications of their one-dimensional fibres,
\[
 R^qg_*\mathcal E=0\ (q\ne2),\qquad
 R^2g_*\mathcal E=\bigoplus_{i=1}^{12}\C_{t_i}.        \tag{B.4}
\]
In particular $(h^0,h^1,h^2,h^3)(Y,\mathcal E)=(0,0,12,0)$.
\end{statement}

\begin{proof}
The transition rule realizes (B.3) as a cubic in a projective plane bundle over $\PP^1$, with Weierstrass coefficients in $\OO(4)$ and $\OO(6)$; hence its total space is projective. Its discriminant is $-16D(t)$. Exact polynomial division gives
\[
 \gcd(D,D')=1,\qquad \gcd(D,t^4+1)=1.
\]
The leading coefficient $31$ is its nonzero value in the chart at infinity. Each discriminant root therefore has nonzero $c_4=-48(t^4+1)$ and is a nodal cubic. A simple discriminant zero is transverse to the nodal discriminant and makes the surface smooth at that node. Alternatively the Groebner basis of the affine total-space singularity ideal $(F,F_x,F_y,F_t)$ is $(1)$; the exact computation is attached. At the zero section smoothness follows from $F_z=1$ in the chart $y=1$. At infinity the coefficients are $a=1+s^4,b=1$, with discriminant $-16\cdot31$ at $s=0$. These checks cover the surface.

The sections $P,O$ are disjoint and meet the smooth locus of every fibre. For $P$, a singular point would require both $a(t)=0$ and $t=0$, which is impossible. Each fibre is an integral arithmetic-genus-one curve with trivial dualizing sheaf. The line bundle $M_E$ has degree zero on every fibre and is nontrivial there. On a smooth fibre, a trivialization would give a meromorphic function with one simple pole and one simple zero, hence a degree-one map to $\PP^1$, contradicting genus one. On a nodal fibre the normalization calculation in (B.2) says that two distinct smooth points determine a nontrivial degree-zero line bundle. Thus $H^q(E_t,M_E^{\pm1}|_{E_t})=0$ for both $q=0,1$: a nonzero section of a degree-zero line bundle on an integral curve is a trivialization, and Serre duality gives the assertion in degree one. Grauert's theorem, applied to the proper flat map $h$, yields
\[
 Rh_*M_E=Rh_*M_E^{-1}=0.                               \tag{B.5}
\]

We compute the cotangent direct image, including its scheme length. Let $Z\subset E$ be the reduced set of the twelve critical points. The relative dualizing sheaf is $\omega_{E/\PP^1}=h^*\OO(1)$. Indeed the residue frame $dx/(2y)$ has this transition under $x_\infty=s^2x,y_\infty=s^3y$, and extends as a dualizing frame at a node. It follows that $K_E=h^*\OO(-1)$. There is an exact sequence of coherent sheaves on $E$
\[
0\longrightarrow h^*\Omega^1_{\PP^1}
 \xrightarrow{dh}\Omega^1_E
 \longrightarrow\mathcal I_Z\otimes\omega_{E/\PP^1}
 \longrightarrow0.                                   \tag{B.6}
\]
Away from $Z$ this is the ordinary relative cotangent sequence. At a node choose coordinates $vw=t$. The map from $\Omega^1_E$ to the relative dualizing line sends $dv\mapsto v\omega$, $dw\mapsto-w\omega$; its image is $(v,w)\omega$ and its kernel is generated by $w\,dv+v\,dw=dt$. This proves (B.6) on every stalk, with the critical scheme reduced.

Tensor (B.6) by $M_E^{-1}$. By (B.5), projection formula, and the sequence for the ideal of $Z$, its derived direct image has only degree one, namely
\[
 R^1h_*(\Omega^1_E\otimes M_E^{-1})
 =h_*(\omega_{E/\PP^1}\otimes M_E^{-1}|_Z),
 \qquad h_*(\Omega^1_E\otimes M_E^{-1})=0.             \tag{B.7}
\]
The right side is the sum of twelve reduced skyscrapers. In particular (B.7) calculates the length rather than merely the residue dimension at each point.

Contraction on the smooth surface gives the canonical isomorphism $T_E\otimes K_E\simeq\Omega^1_E$. Consequently the product tangent decomposition identifies
\[
\mathcal E=(\Omega^1_E\otimes M_E^{-1})\boxtimes\OO(-2)
 \ \oplus\ (K_E\otimes M_E^{-1})\boxtimes\OO.
\]
The second summand has zero derived direct image by (B.5). The first has only the tensor product of degree-one terms: (B.7) and $H^1(\PP^1,\OO(-2))=\C$. The relative coherent Kunneth comparison gives (B.4). The Leray spectral sequence over a curve then gives the four stated global dimensions.

On a singular fibre $S_i=E_{t_i}\times\PP^1$, the restriction of $M_E$ is a nontrivial $A_\lambda$. Statement B1 gives the two vanishings and ambient dimensions $(1,1,0)$. Adjunction for the Cartier fibre gives $\omega_{S_i}=K_Y|_{S_i}$, because its normal line is trivial. Serre duality for this compact Gorenstein surface identifies, for $q=0,1,2$,
\[
 H^q(S_i,\mathcal E|_{S_i})^\vee
 =H^{2-q}(S_i,\Omega_Y^1|_{S_i}\otimes A).
\]
This proves the fibre cohomology. It also explains the fibre value $h^1=1$ while $R^1g_*\mathcal E=0$: tensoring a reduced skyscraper in degree two with the residue field has a degree-one Tor term.
\end{proof}


% B3 from part_B.tex

\begin{statement}{B3}{Proved}{The hypotheses satisfied by the compact example}
The threefold $Y$ of Statement B2 has $H^1(Y,\Z)=0$, $H^2(Y,\Z)\simeq\Z^{11}$, algebraic dimension $3$, and $e(Y)=24$. Its map $g$ is holomorphic, proper, surjective, and has connected fibres. Thus it tests the fibrewise inference in Section 7 of \cite{CDP20}; it does not satisfy the vanishing of $H^2(Y,\Z)$ or algebraic dimension one required by that paper's Theorem 2.2.
\end{statement}
\begin{proof}
The residue calculation above gives $K_E=h^*\OO(-1)$. Relative duality gives $R^1h_*\OO_E=\OO(-1)$, since every fibre is an integral genus-one curve, and $h_*\OO_E=\OO$. Thus $q(E)=p_g(E)=0$ and $H^0(E,K_E^{\otimes m})=0$ for every positive $m$. Castelnuovo's rationality criterion, in the form $q=P_2=0$ for a smooth projective surface, makes $E$ rational. A rational smooth projective surface is obtained from a minimal rational surface by point blowups; its first integral homology vanishes and its second integral cohomology is free. Euler additivity gives $e(E)=12$, since the twelve nodal fibres each have Euler characteristic one and the smooth elliptic fibre has Euler characteristic zero. Poincare duality gives $b_2(E)=10$. The integral Kunneth isomorphism for $Y=E\times\PP^1$ gives the asserted groups and Euler number. Projectivity and dimension give $a(Y)=3$. Connectedness and surjectivity follow directly from the integral cubic fibres of $h$ and the second product factor.
\end{proof}

The construction in Statement B2 supplies the compact local counterexample, including its bundles, cohomology groups and direct-image lengths. The universal assertion under every hypothesis of CDP20 Theorem 2.2 is proved in C4e below: positive third Chern number forces one fixed irreducible nonnormal reduced support whose normalization has the required nonzero group for every ambient Picard twist. Its proof uses the full-fibre comparisons of A.3 and the conductor arguments of A.5. The named-family tests are C4g.


\subsection{B.3. Necessary singular fibres for positive Euler characteristic}

% C4b from part_C4.tex

\begin{statement}{C4b}{Proved}{An unavoidable failure on a rational homology six-sphere}
Let $f:Y\to\PP^1$ be a connected-fibre holomorphic map from a compact complex threefold with $b_1=b_2=b_3=0$. For every $L\in\Pic(Y)$ there is a scheme-theoretic fibre $F$ with
\[
 H^2(F,(T_Y\otimes L)|_F)\ne0,\qquad
 H^0(F,\Omega_Y^1|_F\otimes(K_Y\otimes L^{-1})|_F)\ne0. \tag{C.16}
\]
For torus general fibre the reduced fibre Euler numbers sum to two. Within the fibre class of C4a, exactly one fibre has rational normalization; it is an irreducible reduced normal-crossing fibre of Euler number two. For that normalization,
\[
H^2(\widetilde S,K_{\widetilde S}\otimes
             \eta^*(K_Y^{-1}\otimes L))=\C
\quad\hbox{for every }L\in\Pic(Y).                    \tag{C.17}
\]
\end{statement}
\begin{proof}
The divisor of any fibre is nonzero and effective. The index argument C.1 therefore gives $H^2(Y,T_Y\otimes M)\ne0$ for every line bundle $M$. Suppose $R^2f_*(T_Y\otimes L)=0$. Twisting by a sufficiently large positive pullback from the base kills $H^1(R^1f_*)$, so Leray would give a line bundle $M=L\otimes f^*\OO(n)$ with $H^2(Y,T_Y\otimes M)=0$, a contradiction. Thus $R^2f_*(T_Y\otimes L)$ is nonzero. Its stalk has nonzero residue at some point, and top-degree base change gives the first nonvanishing in (C.16). Cartier-fibre adjunction and Serre duality give the second, with the exact line bundle displayed there.

For torus general fibre C2a gives $\sum e(F_{\rm red})=e(Y)=2$. In the class of C4a, A.4.6 makes every singular reduced contribution either zero or two, with two exactly for the rational hexagonal normalization; smooth finite torus quotients have zero Euler number. There is consequently exactly one positive fibre. Its smooth rational normalization has $\Pic^0=0$ and torsion-free integral second cohomology. Every ambient line bundle has torsion first Chern class, whose pullback is therefore zero. In particular both $\eta^*K_Y$ and $\eta^*L$ are trivial. Serre duality on the connected rational normalization gives (C.17).
\end{proof}


% C4e from part_C4.tex

\begin{statement}{C4e}{Proved}{The unrestricted normalization target under the integral CDP hypotheses}
Let $f:Y\to\PP^1$ satisfy the integral CDP hypotheses and $c_3(Y)>0$. Then a fixed irreducible nonnormal reduced support $S$ has normalization $\eta:\widetilde S\to S$ satisfying
\[
 H^2\bigl(\widetilde S,\omega_{\widetilde S}\otimes
           \eta^*(K_Y^{-1}\otimes L)\bigr)\ne0
 \qquad\hbox{for every }L\in\Pic(Y).                  \tag{C.21}
\]
The assertion concerns the reduced support; no multiplicity-one conclusion is included.
\end{statement}
\begin{proof}
C.3.1 and the exclusions C.3.7, C.3.9 and C.3.10 show
that the general fibre is a two-torus, a Hopf surface, or
an Inoue surface. In the last two cases it has
$h^1(\OO)=1$ and $h^2(\OO)=0$.
For torus general fibre, C4c gives a fixed full fibre with a nonzero ambient cotangent section for every ambient twist. In the second case, C4d gives persistence along $\C\beta$, and A.5.13 identifies this with $\Pic(Y)$. From this point the proof uses precisely that persistence property and does not use $c_3(Y)>0$ again. A.3.5 makes its support irreducible, so it is $mS$ for a single positive integer $m$. A.3.4 excludes a smooth support, and A.5.2 then excludes a normal support. Its filtration in A.3.1 has finitely many graded factors
\[
 \Omega_Y^1|_S\otimes M|_S\otimes\OO_Y(-jS)|_S,
 \qquad 0\le j<m.
\]
Nonvanishing on the filtered sheaf requires nonvanishing on at least one factor. Each factor's nonvanishing locus is closed analytic in the connected vector group $\Pic(Y)$. A finite union of proper such loci cannot cover that group. Hence one factor has nonzero sections for every $M$. Translating by $\OO_Y(-jS)$ shows
$H^0(S,\Omega_Y^1|_S\otimes M|_S)\ne0$ for every $M$.
A.4.9 kills the torsion-free quotient on a dense open set of twists. The exact sequence A.1.1 and proper semicontinuity then show that the inverse-Jacobian kernel itself has sections for every twist.

We apply the precise persistent-kernel reduction A.5.4. It gives a fixed component of the normalization on whose resolution every twist of one fixed rank-one sheaf has a section. The Gauduchon degree argument there proves that the pullback of every ambient line has degree zero for every Gauduchon metric. Its conductor comparison excludes nonnegative Kodaira dimension, and its class VII argument makes every ambient line trivial on that normalization component. Rational components have the same triviality because their Picard group has no connected part. Either case gives (C.21) immediately by Serre duality and descent from resolution. It remains to treat an irrational ruled resolution $p:Z\to B$.

Let $h:Z\to S\subset Y$. A.5.3 supplies an effective conductor divisor $D_Z$ with
\[
 K_Z+D_Z=h^*(K_Y+S),\qquad D_Z\cdot F=2              \tag{C.22}
\]
for a general ruling fibre $F$. Its horizontal part is therefore either reduced of total ruling degree two, or $2C$ for a section $C$. We record the comparison through each ruling blowdown $\sigma:Z\to Z'$ used below. Put $D'=\sigma_*D_Z$, let $a$ be the coefficient of its exceptional curve $E$ in $D_Z$, and put $d=\operatorname{mult}_{\sigma(E)}D'$. Then
\[
 K_Z+D_Z=\sigma^*(K_{Z'}+D')+(1+a-d)E.
 \tag{C.22b}
\]
Intersection with $E$, using $c_1(K_Z+D_Z)=0$ and $E^2=-1$, gives $a=d-1$. Thus the log-canonical divisor is the exact pullback of $K_{Z'}+D'$, and the latter again has zero rational Chern class. No factorization of $h$ through the blowdown is asserted or needed.

We first treat the reduced horizontal part $H$. If a component is exceptional for the normalization resolution, every ambient line restricts trivially to it. Its normalization maps finitely onto $B$, and the norm composed with pullback on $\Pic^0(B)$ is multiplication by the covering degree. Hence this restriction has finite kernel; the connected ambient Picard image is trivial, giving (C.21). We may therefore assume that the horizontal components map to conductor curves. Adjunction on each of them gives
\[
 2p_a(H_i)-2+(D_Z-H_i)\cdot H_i=0.
\]
The second term is nonnegative. Since $H_i$ dominates a positive-genus base, its arithmetic genus is at least one. Thus the base has genus one, each $H_i$ is smooth elliptic, and the horizontal curves are disjoint from each other and from the vertical part. Contract ruling $(-1)$-curves to a minimal bundle. A contracted curve outside $D_Z$ has $D_Z\cdot E=1$, so cannot identify two horizontal points or create a horizontal node. A contracted curve in the vertical part is disjoint from $H$. Downstairs the horizontal curves remain smooth and disjoint, and any nonzero vertical divisor on a minimal ruled bundle meets them. Consequently the downstairs vertical part is zero. Reversing the contractions, a blowup away from the horizontal conductor would give coefficient $-1$ on its exceptional curve in (C.22), which persists under further strict transform and contradicts effectiveness. A blowup at a smooth point of a coefficient-one horizontal curve gives exceptional coefficient zero. Hence $D_Z=H$ has no vertical component. Every codimension-one singularity of $S$ is then an ordinary node, and A.2.3 identifies the inverse-Jacobian kernel with the normalization kernel globally, including its codimension-two points. A section of a topologically trivial line on the Kahler resolution has an effective zero divisor of zero Kahler volume, so is nowhere zero. Persistence for all twists therefore makes all ambient pullbacks trivial and proves (C.21).

We now have $D_Z=2C+V$, where $C$ is a section and $V$ is vertical. If $C$ is exceptional for $Z\to\widetilde S$, every ambient line is trivial on $C$. Restriction $\Pic^0(Z)\to\Pic^0(C)$ is an isomorphism, because $C$ is a ruling section. This again makes all ambient pullbacks trivial. Otherwise let $D\subset S$ be its image and let $D^\nu$ normalize that curve. When $D^\nu$ is rational, every ambient line pulls back trivially to $D^\nu$, hence to $C$, giving the same conclusion.

Suppose instead that $g(D^\nu)\ge1$. A vertical rational curve cannot map finitely onto $D^\nu$, so all the normalization branches over the general point of $D$ belong to $C$. Their conductor coefficient is two. We derive the possible transverse reduced plane curves rather than assuming a list of singularities. Write its equation as $f=\prod_i f_i$ in the smooth-plane local ring and let $R_i$ and $B_i$ be a branch ring and its normalization. The elements of $R=\OO/(f)$ supported only on branch $i$ form
\[
 \left(\prod_{j\ne i}f_j\right)\bigg|_{f_i}\,R_i.
\]
This equality follows from
$\bigcap_{j\ne i}(f_j)=(\prod_{j\ne i}f_j)$ in that unique
factorization domain. Multiplying by arbitrary elements of
$B_i$ therefore shows that the conductor exponent on branch $i$
is its intrinsic exponent plus $\sum_{j\ne i}I(f_i,f_j)$.
A singular branch has intrinsic exponent at least two:
exponent one would imply $\C+t\C\{t\}=\C\{t\}\subset R_i$.
Exponent two forces $R_i=\C+t^2\C\{t\}=\C\{t^2,t^3\}$,
since an element of order one would again make the ring smooth.
Any additional branch has intersection multiplicity at least
two with a singular branch and is therefore excluded. Thus a
singular branch gives exactly the cusp. When all branches are
smooth, their pairwise positive intersection multiplicities sum
to two on each branch. This gives two branches of contact two
(the tacnode), or three pairwise transverse branches (the
ordinary triple point), and exhausts the possibilities.
The only remaining alternative is a nonconstant pullback
$\Pic(Y)\to\Pic(Z)$; a constant pullback is the zero
homomorphism and already gives (C.21). In that alternative,
the coherent conductor argument A.5.12 gives $p_a(D)\le1$
for torus general fibre, and A.5.13 gives the same bound
in the second case.
Since $g(D^\nu)\ge1$, it follows that $D$ itself is smooth
elliptic. This conclusion permits nonreduced conductor image
schemes and imposes no assumption on the original singular
curve. We can therefore apply C2d directly. It makes the
transverse type constant, gives the etale conductor cover
$C\to D$, and proves $C^2=0$. In particular $C$ and the
ruling base are elliptic.

Adjunction in (C.22) gives $V\cdot C=0$. No ruling $(-1)$-curve meets $C$: for one outside $V$, its intersection with $K_Z+2C+V$ would be at least $-1+2=1$; for one in $V$, the equality $V\cdot C=0$ forbids such an intersection. Thus every contraction to the minimal ruled bundle avoids $C$ and preserves its self-intersection zero. On that bundle $K+2C$ is numerically trivial, so the effective vertical part is numerically trivial and hence zero. A first reverse blowup, necessarily away from $C$, would again produce negative conductor coefficient. There are therefore no reverse blowups: $Z$ is minimal ruled and $D_Z=2C$. There is also no exceptional curve for the normalization resolution: a negative curve $E\ne C$ would satisfy $2p_a(E)-2=E^2-2C\cdot E<0$, hence would be rational and contained in a ruling fibre; a minimal ruled bundle has only self-intersection-zero fibres. Thus the normalization itself is smooth here.

By the explicit inverse-Jacobian computations in C2d, its pullback allows at most a simple pole along $C$. There are no other divisorial poles. Hartogs extension for the line bundle on the smooth surface gives an injection of the kernel into $h^*\OO_Y(-S)\otimes\OO_Z(C)$. The exact sequence of the section $C$ and $R^1p_*\OO_Z=0$ give
\[
 0\longrightarrow\OO_B\longrightarrow p_*\OO_Z(C)
   \longrightarrow N_{C/Z}\longrightarrow0,
 \qquad \deg N_{C/Z}=0.                              \tag{C.22a}
\]
For a degree-zero line $P$ on the elliptic curve, both end terms after tensoring by $P$ have zero sections unless $P=\OO_B$ or $P=N_{C/Z}^{-1}$. Thus $H^0(Z,\OO_Z(C)\otimes p^*P)$ is supported at at most two points of $\Pic^0(B)$. The ambient pullback image is a connected complex subgroup of this Picard group. Persistence for every ambient twist forces that subgroup to be a point, and it contains the identity. Therefore every ambient line pulls back trivially to $Z$ in this last case as well.

The equality $\pi_*\OO_Z=\OO_{\widetilde S}$ for the normal resolution descends triviality of a pulled-back line: projection formula identifies its space of sections with that of its trivial pullback, and the same statement for its inverse gives mutually inverse trivializations. Consequently one normalization component has trivial $\eta^*(K_Y^{-1}\otimes L)$ for every $L$. Serre duality on that compact normal Cohen--Macaulay surface identifies its top dualizing cohomology with $H^0(\OO)^\vee=\C$. It is a direct summand of the target in (C.21), proving the claim.
\end{proof}

The local example A.2.5 has actual vertical inverse-Jacobian
poles, and A.2.6 has a critical curve that is not a local
complete intersection. Their local rings remain part of the
audit. The global conductor-duality arguments A.5.12--A.5.13 prove
that a positive-genus conductor image in the persistent ruled
case is smooth, and the exact contractions above then remove
the vertical conductor. Thus neither local phenomenon leaves
an unproved step in the stated theorem. The assertion
is about the normalization of a reduced support and does not
claim that its full fibre has multiplicity one.


% C4g from part_C4.tex

\begin{statement}{C4g}{Proved}{The necessary normalization defect in the three named test families}
For every admissible member of the original $(3,4)$ family,
every member of the two-parameter family in E13, and every
member of the Liu--Rollenske family in E14, there is a reduced
full hexagonal fibre $W$ with normalization
$\eta:\mathrm dP_6\to W$ such that
\[
 \eta^*L=\OO_{\mathrm dP_6}\quad(L\in\Pic(Y)),\qquad
 H^2\bigl(\mathrm dP_6,K_{\mathrm dP_6}\otimes
                    \eta^*(K_Y^{-1}\otimes L)\bigr)=\C.
 \tag{C.TEST.1}
\]
In these three families the singular-fibre Euler contributions are
\[
\begin{array}{c|c|c}
\text{family}&\text{reduced singular supports}&
                 \text{their Euler numbers}\\ \hline
(3,4)&W&2\\
\text{two-parameter E13}&W,\ W_1,\ W_2&2,0,0\\
\text{Liu--Rollenske E14}&W,\ C_{\rm cusp}&2,0 .
\end{array}
\tag{C.TEST.2}
\]
The other critical full fibres are smooth multiple fibres,
with multiplicities $(3,4)$, $4$, and $4$, respectively.
Every one of those full multiple fibres passes the generic
ambient vanishing test A.3.4, including all its nilpotent
layers. Consequently positive $c_3$ in these examples does
not require a failure of that test on a nonreduced fibre.
The reduced full fibre $W$ already supplies the fixed
normalization obstruction, and $H^2(Y,T_Y\otimes L)\ne0$
for every ambient line bundle in each family.
\end{statement}
\begin{proof}
The original local filling is the hexagonal toric quotient
in A.4.1. The two-parameter construction in E13 places this
same filling at the simple nodal value where its specified
linearization vanishes; its other two simple nodal values
have rank-one fillings. The coordinate and line-bundle
comparisons of E14 place the hexagonal filling at its
separate simple node and the cuspidal quotient at the
remaining critical value. Those proofs construct the actual
local spaces and their identifications over the punctured
base discs. Thus the normalization of $W$ in each case is
the stated smooth rational surface.

These threefolds have $H^2(Y,\Z)=0$. Pulling back any
ambient line to $\mathrm dP_6$ gives zero integral first
Chern class. The exponential sequence there, with
$H^1(\mathrm dP_6,\OO)=0$, makes that line trivial.
This applies to $K_Y^{-1}\otimes L$ as well as to $L$.
Serre duality on $\mathrm dP_6$ proves (C.TEST.1).
The last assertion about ambient $H^2$ follows from C4f
through the zero Picard pullback at this fixed fibre.
It does not require surjectivity of the global
normalization map for every twist; C5b computes the
actual exceptions to that surjectivity in the original
family.

The toric strata of $W$ consist of one $(\C^*)^2$ stratum,
three $\C^*$ strata and two triple points. Their compactly
supported Euler numbers give $e(W)=2$.
Each rank-one normalization is a $\PP^1$-bundle over an
elliptic curve, with its zero and infinity sections glued.
Both the normalization and these conductor curves have
Euler number zero, so normalization additivity gives
$e(W_i)=0$. At the cusp in E14, normalization of the
cuspidal rational curve makes its degree-zero line bundle
trivial. The lifted contracting action on this normalized
line complement is a product of a fixed automorphism of
$\PP^1$ and multiplication by a constant of modulus
strictly less than one on $\C^*$. Its quotient is a
$\PP^1$-bundle over the elliptic curve given by that
multiplicative period. The cusp normalization is
bijective; a proper continuous bijection between these
compact Hausdorff spaces is a homeomorphism. Hence
$e(C_{\rm cusp})=0$. This proves (C.TEST.2).

Finally the explicit good-reduction covers in E13--E14
and the original root charts have the stated orders.
Their reduced central fibres are smooth finite unramified
torus quotients, and their Euler numbers are zero.
A.3.4 applies to their actual full fibres, proving the
asserted generic ambient vanishing. Euler localization
therefore gives $c_3(Y)=2$ in all three rows, while the
nonreduced fibres satisfy the complete multiple-fibre
test. This is the required test of the proposed additional
nonreduced-fibre necessity.
\end{proof}

\paragraph{Literature comparison for the geometric replacement.}
The ambient-kernel failure for the original construction is
located in \cite[Section~10]{AC26}; the general kernel is
calculated in A.1--A.2. Statements C4e--C4f use the actual
conductor-duality map and the independently proved
smooth-fibre exclusions to obtain a geometric criterion
under all the integral hypotheses of \cite[Theorem~2.2]{CDP20}.
The latter paper already states the Kodaira and bielliptic
exclusions in its Corollary~5.3, with its Corollary~5.2 as
an input; C.3.7 and C.3.9--C.3.10 prove those exclusions
without that input. The Inoue exclusion C.3.12 is derived
from the explicit complete families of \cite{Inoue1974}.
Searches on 15 November 2026 for compact threefolds,
vanishing second Betti number, Kodaira or Inoue general
fibres, normalization obstructions, and degenerate-torus
Euler numbers did not supply the all-singularity
criterion C4f in the sources opened. We claim no priority.

